% =======================================================================================
\section{Method}
\label{sec:method}
% =======================================================================================

\subsection{Flow Coordinates}
\label{sec:method-chart}

The flow transports points of $\R^{8S}$, one chart image per source (\cref{eq:chart}). Frequency
is scaled to the window, $\log\fdot$ standardised given $f_0$, and the sky taken in galactic
coordinates scaled by the prior's widths. Amplitude, inclination, polarisation and phase enter
as the two helicity amplitudes
$\zeta_\pm\propto\Amp\,(1\mp\cos\iota)^2\,e^{\ii(\phi_0\pm2\psi)}/4$, in units of the noise, with
their radius warped by $\operatorname{asinh}$. The signal is linear in $\zeta_\pm$
(\cref{eq:forward-model-helicity}), so the four copies of a source in $(\psi,\phi_0)$ map to one
point, and at high \ac{snr} the likelihood in $\zeta_\pm$ is close to Gaussian
(\cref{app:chart-helicity}). The chart is flat, so the flow's paths are straight lines
(\cref{app:chart-flat}). The sources are slots in order of $f_0$, so each mode of the posterior
has one copy instead of $S!$, one per labelling (\cref{app:chart-slots}).

\subsection{Flow Matching}
\label{sec:method-fm}

The network first reads a Gaussian frame off the data, a mean $\mu$ and a standard deviation
$\sigma$ per coordinate of each slot, around the prior's moments (\cref{eq:fm-frame}). It is trained
on the Gaussian likelihood of the true sources $X_1$, so it learns the posterior's marginal mean and
width, with a weight that keeps loud sources, whose frames are decades narrower than the prior, from
dominating the batch (\cref{eq:fm-frame-loss}). The flow then works in that frame, $Z =
(X-\mu)/\sigma$: paths are straight, $Z_s = Z_0 + s(Z_1-Z_0)$, from a standard normal $Z_0$ to the
true sources $Z_1$ (\cref{eq:fm-path}), and the flow cannot move the frame. The network predicts
the mean $\hat Z_1$ and log standard deviation $\log\tau$ of a Gaussian over the endpoint, per
coordinate, and is trained on the endpoint's negative log-likelihood (\cref{eq:loss}) at times
uniform along the path; its velocity is $(\hat Z_1 - Z_s)/(1-s)$ (\cref{app:fm-loss}). A loud
source's posterior is a small fraction of a bin wide in $f_0$. A Gaussian model
(\cref{prop:fm-sensitivity}, \cref{fig:fm-sensitivity}) shows that the velocity loss of standard
flow matching barely rewards narrowing such a posterior, while the likelihood pays the same for
every e-fold of excess width, relative to each source's own width. The same model shows that a
uniform sampler grid run on the prior's scale collapses the narrowest posteriors, and that in a
frame of the posterior's width it does not; draws use \FlowMatchingOdeSteps{} midpoint steps of the
velocity and are mapped back with $X = \mu+\sigma Z$ (\cref{app:fm-sampling}).

\subsection{Network}
\label{sec:method-network}

The network is a two-stream transformer \citep{peebles2023scalable,esser2024scaling} with a data
stream and a source stream (\cref{fig:network-architecture}). The data stream reads the window as
a short-time Fourier transform, \NetworkPatch{} bins per token, each token holding its band's time
samples in the three channels (\cref{eq:network-stft}): \NetworkTokensNarrow{} tokens in the
narrow band, \NetworkTokensFull{} in the full band. One frame register per slot runs with the
data stream and gives that slot's frame (\cref{app:network-frame}). The source stream holds one
token per slot, embedded from its coordinates in the frame and in the chart and from the frame
itself, together with attention sinks \citep{xiao2024efficient,darcet2024vision}. Attention is
one-way: the data stream and the registers attend only to each other and are conditioned on the
window alone, never on the time or the sources, so their keys and values, and the frame, are
computed once per window and reused by every draw and every sampler step; the source stream
attends to everything and is conditioned on the time. Rotary embeddings in frequency
\citep{su2024roformer} place each source at its $f_0$ and each data token at its centre bin.
\Cref{app:network} gives the embeddings and \cref{tab:network-sizes} the sizes; all results use
size B, \num{\NetworkBTotalFull} parameters.

\subsection{Training}
\label{sec:method-training}

Every batch is simulated on the device: \TrainingObservations{} windows, each with
\TrainingDraws{} draws of time and base point sharing its cached data stream, so no observation is
seen twice (\cref{app:training-batch}). The optimiser is schedule-free
\citep{defazio2024road} around Muon \citep{jordan2024muon,li2025normuon} for the hidden weight
matrices and AdamW for the rest, and evaluation uses the averaged weights
(\cref{app:training}). The narrow band trains for \num{\TrainingStepsNarrow} steps and the full
band for \num{\TrainingStepsFull}, \qty{\TrainingComputeNarrowBHours}{\hour} and
\qty{\TrainingComputeFullBHours}{\hour} on one A100 (\cref{app:training-compute}).
